% ==============================================================================
% Plantilla Institucional de Trabajo Terminal — UPIIZ IPN
% Archivo: capitulos/01-introduccion.tex
% ==============================================================================

\chapter{Introducción}
\label{cap:introduccion}

% ==============================================================================
% REGLA INSTITUCIONAL:
% A partir de esta sección inicia la numeración de páginas en números arábigos (1, 2, 3...).
%
% INSTRUCCIONES PEDAGÓGICAS:
% La introducción debe contextualizar el proyecto, describir la motivación y el
% alcance general del trabajo mecatrónico desarrollado, así como presentar un
% resumen estructurado del contenido de los capítulos que componen el documento.
% ==============================================================================

% Escriba aquí la descripción general y contextualización de su proyecto...
El presente Trabajo Terminal aborda el diseño y desarrollo de una solución mecatrónica orientada a resolver las necesidades identificadas en el ámbito de la ingeniería aplicada. En este documento se presentan de manera secuencial los fundamentos teóricos, metodológicos y de ingeniería que sustentan el proyecto.

A continuación se detalla la estructura capitular del reporte:
\begin{itemize}
    \item El \capref{cap:justificacion} expone la justificación técnica, social y económica del proyecto.
    \item El \capref{cap:antecedentes} presenta los antecedentes y trabajos previos relacionados.
    \item El \capref{cap:marco-teorico} sintetiza los fundamentos teóricos y modelos físico-matemáticos empleados.
    \item El \capref{cap:estado-del-arte} revisa el estado del arte y las tecnologías vigentes afines.
    \item El \capref{cap:planteamiento} define el planteamiento del problema y la solución propuesta.
    \item El \capref{cap:desarrollo} describe el desarrollo técnico detallado de la ingeniería mecatrónica.
    \item El \capref{cap:validacion} presenta el análisis y validación conforme a las especificaciones planteadas.
    \item El \capref{cap:conclusiones} establece las conclusiones derivadas del cumplimiento de los objetivos.
    \ifTTII
    \item El \capref{cap:trabajo-futuro} plantea las líneas de trabajo a futuro e iteraciones posteriores.
    \fi
\end{itemize}

% ==============================================================================
% OBJETIVOS DEL PROYECTO
% Copie aquí los objetivos EXACTAMENTE como fueron aprobados en el protocolo registrado.
% NO los modifique ni parafrasee sin aprobación formal del comité evaluador.
% ==============================================================================

\section{Objetivo general}
\label{sec:obj-general}

% Escriba aquí el objetivo general aprobado en el protocolo...
Desarrollar un sistema mecatrónico integral para la automatización y control del proceso definido, mediante el diseño e integración de subsistemas mecánicos, electrónicos y de software, a fin de validar su rendimiento ante especificaciones operativas cuantitativas.

\section{Objetivos específicos}
\label{sec:obj-especificos}

% Escriba aquí la lista de objetivos específicos aprobados en el protocolo...
\begin{enumerate}
    \item Establecer la lista de requerimientos y especificaciones de diseño mecatrónico a partir de las condiciones de operación.
    \item Diseñar y modelar los componentes mecánicos, estructurales y cinemáticos requeridos por el sistema.
    \item Diseñar y seleccionar la instrumentación electrónica, sensores, actuadores y etapas de potencia.
    \item Desarrollar los algoritmos de control, lógica de programación e interfaz de usuario para el sistema embebido.
    \item Evaluar y validar cuantitativamente el desempeño global del sistema mediante pruebas experimentales y análisis de error.
\end{enumerate}
